\chapter*{Abstract}
\addcontentsline{toc}{chapter}{Abstract}
\markboth{Abstract}{}
\begingroup\setstretch{1.2}

\paragraph{Background.}
Diabetic retinopathy (DR) affects an estimated 103 million adults worldwide and is a
leading cause of preventable blindness in the working-age population.
Automated screening is clinically essential, yet no system identified in the reviewed
literature resolves three co-existing structural limitations at once: L1 (no relational
context between patients), L2 (no global vascular topology), and L3 (no clinical safety
mechanism: every image is graded, and the Severe-DR Non-Referral Rate, Sev-NR, stays above
the $2\%$ threshold).

\paragraph{Methods.}
This thesis is an eleven-publication research programme. Three hypotheses are pre-specified,
each with its test and rejection threshold: H1, an inductive population graph adds
information that a per-image classifier cannot recover; H2, $H_0/H_1$ persistent homology
of the retinal vasculature is associated with severity; H3, the combination of the three
axes is sufficient (H3a) and, within the evaluated architectural family, necessary (H3b)
for seven predefined computational criteria. Foundations established on distributed
systems are transferred to DR grading
through hybrid models, the
TOPORET population-graph framework
and the DOTS system.

\paragraph{Results.}
H1 is confirmed ($t_4=3.82$, one-tailed $p=0.009$; $+4.0$ percentage points of quadratic
weighted kappa, QWK). H2 is confirmed ($F(4,4995)=1308.9$, $p<0.001$, $\eta^2=0.51$).
On five public benchmarks ($N>108{,}000$ images; $10{,}000$ held-out test predictions),
DOTS reaches QWK $=0.951$, AUC $=0.988$, ECE $=0.040$, Sev-NR $=1.9\%$, coverage
$93.8\%$ and $0\%$ ordinal rank violations: H3a is met at the point estimate, but the
95\% confidence interval of Sev-NR, $[1.4\%,\,2.4\%]$, still crosses the $2\%$ line.
H3b holds by exhaustive enumeration: all six proper subsets of the three axes fail
Sev-NR $<2\%$.

\paragraph{Conclusions.}
The three limitations can be resolved in one architecture, and no proper subset of it
meets all predefined criteria. This is Level~1 (computational, retrospective) evidence,
not a licence for clinical use; a four-phase pathway (2026--2029) towards prospective
validation and CE/FDA submission is proposed. The three-axis framework is offered as a
template for any graded disease with an ordinal scale, a structural anatomical prior and
an international safety threshold.

\vspace{0.6em}
\noindent\textbf{Keywords:} diabetic retinopathy; graph neural networks; GraphSAGE;
persistent homology; ordinal regression; CORAL; selective prediction; clinical AI safety;
RETFound; medical device software.
\par\endgroup
